\section{Problem Formulation and Method}
\label{sec:method}

\subsection{Overview: Seven-Step Pipeline and Dual Role of the Expert}
\label{sec:method_overview}
Our methodology spans two tightly linked parts---\textbf{Part~I (Steps~1--5)}: a controlled benchmark of low-dimensional intent interfaces under channel degradation, and \textbf{Part~II (Steps~6--7)}: deploying the winning goal-level interface in an end-to-end VLM manipulation system:
\begin{enumerate}[leftmargin=1.5em,itemsep=1.5pt,topsep=2pt]
  \item \textbf{Step~1 (Reference expert training \& freezing):} Train a goal-visible PPO policy per task and freeze its final checkpoint to define the 100\% performance ceiling $P_E$.
  \item \textbf{Step~2 (Deterministic rollout \& PCA basis):} Roll out the frozen reference expert to collect 1.23M continuous transitions and fit orthogonal action-synergy PCA bases $W$.
  \item \textbf{Step~3 (Intent generation \& channel degradation):} At each step, compute a $k$-D intent $z_t$ (action, object-subgoal, or goal level) and pass it through controlled zero-order hold, delay, or quantization to obtain $\tilde{z}_t$.
  \item \textbf{Step~4 (Goal-hidden student training):} Strip goal observations from the actor and train a goal-hidden student via asymmetric-critic PPO using $\tilde{z}_t$ as its sole task cue, executing all interface and rate/delay sweeps here.
  \item \textbf{Step~5 (Deterministic evaluation):} Evaluate all 25 periodic checkpoints across 1,024 fixed-seed episodes to compute exact completions, recovery, staleness, and waiting-time fits.
  \item \textbf{Step~6 (Unified 60~Hz multi-skill controller):} Train a single goal-conditioned controller (\texttt{Unified v4}) for the two-skill hinged cube and verify its 96-state BFS execution ceiling and 4~s pause stability.
  \item \textbf{Step~7 (Upstream VLM adaptation \& online 60~Hz loop):} Adapt \texttt{Qwen3-VL-8B-Instruct} via Map-then-Reason LoRA SFT and closed-loop GRPO to drive frozen \texttt{Unified v4} from live 3-camera images.
\end{enumerate}

\paragraph{Clarifying the dual role of the expert.}
To avoid ambiguity, note that the frozen expert serves \emph{two distinct roles} across Parts~I and~II.
In \textbf{Part~I (Steps~1--5)}, it is a \textbf{Reference Expert}: it sets the $P_E$ baseline, supplies the offline PCA basis, and—\emph{only} for action-level intents—acts online as an action-intent oracle computing $a_t^E \to z_t$.
Crucially, student training in Step~4 is \emph{not} teacher--student action distillation; every student learns autonomously from environment rewards via asymmetric-critic PPO.
In \textbf{Part~II (Steps~6--7)}, once Part~I establishes that slow planners must emit goal-level commands, frozen \texttt{Unified v4} serves directly as the \textbf{60~Hz Goal-Conditioned Skill Controller} (analogous to an arm's IK controller) executed under the slow VLM planner.

\subsection{Tasks, Simulation Setup, and Evaluation Protocol}
\label{sec:method_tasks}
All experiments run in Isaac Lab v2.3.1 \citep{mittal2023orbit} on the 24-DOF Shadow Dexterous Hand (20 actuated joints: 2 wrist drives at action indices 0--1 and 18 finger drives at indices 2--19, with 4 coupled distal joints).
Physics runs at 120~Hz with decimation 2, yielding a 60~Hz discrete-time control loop and 10~s episodes (600 control steps) with continuous observations $s_t \in \mathbb{R}^{157}$ and continuous actions $a_t \in [-1,1]^{20}$.
Observations include joint positions/velocities, object pose/velocities, goal quaternion \texttt{obs[64:68]}, relative object--goal quaternion \texttt{obs[68:72]}, fingertip states, and previous action.
We evaluate three tasks of increasing complexity:
\begin{itemize}[leftmargin=1.5em,itemsep=1.5pt,topsep=2pt]
  \item \textbf{Task~1 ($\mathrm{SO}(3)$ Reorientation):} Standard \texttt{Isaac-Repose-Cube-Shadow-Direct-v0}. Reorient a 6~cm cube to random 3D target orientations within 0.1~rad; each success samples a new target without resetting the hand ($P_E=26.11$).
  \item \textbf{Task~2 ($z$-Axis Reorientation, M2):} \texttt{Isaac-Repose-Cube-Shadow-ZAxis-Direct-v0}, restricting targets to yaw rotations about the palm-normal $z$-axis ($P_E=31.14$).
  \item \textbf{Task~3 (Simplified Rubik's-Like Hinged Cube):} A custom two-part 6~cm cube (0.216~kg, friction 1.0) whose top third twists about the cube $z$-axis against a 0.05~N$\cdot$m detent torque cap every $90^\circ$. It combines whole-cube reorientation (\texttt{any}: arbitrary 3D target; \texttt{up}: top layer upright with random yaw) and $\pm 90^\circ$ top-layer twists (\texttt{turn}: hinge within 0.1~rad of target for 10 consecutive steps, 0.17~s).
\end{itemize}

\paragraph{Exact per-episode evaluation (Step~5).}
Because Isaac Lab's built-in \texttt{consecutive\_successes} metric is an across-reset exponential moving average that understates throughput ($\approx 15.5$ vs.\ $26.11$), all results use an independent deterministic evaluator over 1,024 fixed-seed episodes per checkpoint, reporting exact per-episode completions $P$, drop rate, and normalized recovery:
\begin{equation}
  \mathrm{Recovery}(c) = \frac{P(c) - P_{\mathrm{blind}}}{P_E - P_{\mathrm{blind}}},
  \label{eq:recovery}
\end{equation}
where $P_{\mathrm{blind}}$ is the goal-hidden student with no intent ($0.006$ on Task~1; $19.66$ on Task~2).
Policies train with RL-Games PPO \citep{schulman2017proximal,makoviichuk2021rlgames} for 5,000 iterations (2,048 envs, 163.84M transitions; 8,192 envs on H200 for unified experts), with 3--5 seeds per condition and two-sided Welch $t$- / Mann--Whitney $U$ tests.

\subsection{Goal-Hidden Student Training and Intent Taxonomy (Steps 1--4)}
\label{sec:method_student}
In a preliminary goal-visible stage (Stage~1; 21 runs on Task~1 and 30 on Task~2), appending intents to a goal-visible actor yielded no significant speed or AUC gain (one-way ANOVA N80 $p=0.82$, AUC $p=0.65$; Appendix~\ref{app:stage1}), because the actor already observes the goal directly.
To isolate the information carried by the intent channel, we adopt a \emph{goal-hidden asymmetric-critic} protocol (Figure~\ref{fig:protocol}): the 8 goal fields \texttt{obs[64:72]} are stripped from the student actor (leaving 149-D proprioceptive/object state plus the $k$-D intent $\tilde{z}_t$), while the central value critic \citep{pinto2018asymmetric} retains the full 157-D state during training for low-variance advantage estimation.
A goal-visible asymmetric-critic gate run matches the symmetric expert within 1.4\% ($25.76 = 0.986\,P_E$), confirming no architectural confound.

\begin{figure}[t]
  \centering
  % Compact TikZ diagram for Figure 1 (v2): Part I Student Training & Channel-Sweep Architecture (Steps 1--5)
% Explicitly decomposes: 157-D -> Frozen Expert -> 20-D action -> Fixed Offline PCA Projection -> k-D intent,
% explicit Concat(149-D physical state, z_tilde), and solid-line Asymmetric Critic for all training conditions.
\definecolor{navink}{HTML}{17213A}
\definecolor{mutedgray}{HTML}{62708D}
\definecolor{cardbg}{HTML}{F6F8FC}
\definecolor{cardborder}{HTML}{CBD5E1}
\definecolor{intentblue}{HTML}{2C67F2}
\definecolor{intentbluebg}{HTML}{EEF4FF}
\definecolor{chanorange}{HTML}{D9602F}
\definecolor{chanorangebg}{HTML}{FFF3ED}
\definecolor{criticgreen}{HTML}{1E9C77}
\definecolor{criticgreenbg}{HTML}{EDF9F5}
\definecolor{stripcrimson}{HTML}{C62828}

\resizebox{\linewidth}{!}{%
\begin{tikzpicture}[
    font=\sffamily\footnotesize,
    >=Stealth,
    box/.style={
      rectangle, rounded corners=3.5pt, draw=cardborder, fill=cardbg,
      line width=0.85pt, align=center, inner sep=3.5pt, text=navink
    },
    envbox/.style={
      box, draw=navink!70, fill=cardbg, line width=1.0pt,
      minimum width=2.45cm, minimum height=1.85cm
    },
    expertbox/.style={
      box, draw=navink!65, fill=cardbg, line width=0.95pt,
      minimum width=2.7cm, minimum height=0.82cm
    },
    pcabox/.style={
      box, draw=intentblue!85, fill=intentbluebg, line width=0.95pt,
      minimum width=2.9cm, minimum height=0.82cm
    },
    basisbox/.style={
      box, draw=mutedgray, fill=white, dashed, line width=0.8pt,
      minimum height=0.48cm, inner sep=3pt
    },
    chanbox/.style={
      box, draw=chanorange, fill=chanorangebg, line width=1.0pt,
      minimum width=3.6cm, minimum height=0.82cm
    },
    actorbox/.style={
      box, draw=intentblue, fill=intentbluebg, line width=1.05pt,
      minimum width=3.6cm, minimum height=0.95cm
    },
    criticbox/.style={
      box, draw=criticgreen, fill=criticgreenbg, line width=1.0pt,
      minimum width=3.6cm, minimum height=0.88cm
    },
    concatnode/.style={
      circle, draw=intentblue, fill=intentbluebg, line width=1.0pt,
      inner sep=1.2pt, minimum size=0.46cm, font=\sffamily\small\bfseries, text=intentblue
    },
    wire/.style={->, line width=0.9pt, draw=navink!80},
    criticwire/.style={->, line width=0.95pt, draw=criticgreen!90!black},
    dashwire/.style={->, dashed, line width=0.8pt, draw=mutedgray},
    stripwire/.style={-, dashed, line width=0.9pt, draw=stripcrimson}
  ]

  % Main nodes (vertically compacted to save page height)
  \node[envbox] (env) at (-0.35, 1.30) {
    \textbf{Isaac Lab Env}\\[1.5pt]
    {\scriptsize Shadow Hand}\\
    {\scriptsize 60~Hz ctrl (600 steps)}
  };

  % Top row: 157-D -> Frozen Expert -> 20-D action -> Online PCA Proj -> Intent Channel
  \node[expertbox] (expert) at (5.05, 3.15) {
    \textbf{Frozen Expert $\pi_E$}\\[0.5pt]
    {\scriptsize Step 1 (action oracle)}
  };

  \node[pcabox] (pca) at (9.15, 3.15) {
    \textbf{Online PCA Proj.}\\[0.5pt]
    {\scriptsize $z_t = W_k^\top (a_t^E - \mu)$}
  };

  \node[basisbox] (basis) at (9.15, 4.25) {
    \textbf{Offline PCA Basis (Step 2):} {\scriptsize fixed $\mu, W_k$ (1.23M rollout)}
  };

  \node[chanbox] (channel) at (13.95, 3.15) {
    \textbf{Intent Channel (Step 3)}\\[0.5pt]
    {\scriptsize Hold $N$ (60--2~Hz) $\cdot$ Delay $d$ $\cdot$ $b$-bit}
  };

  % Middle row: Concat + Student Actor
  \node[concatnode] (concat) at (10.55, 1.40) {$\oplus$};

  \node[actorbox] (actor) at (13.95, 1.25) {
    \textbf{Student Actor $\pi_\theta$ (Step 4)}\\[0.5pt]
    {\scriptsize Input: $[\text{149-D state},\; \tilde{z}_t]$}\\
    {\scriptsize Trained via PPO (not distillation)}
  };

  % Bottom row: Asymmetric Critic (solid -- required in every training condition)
  \node[criticbox] (critic) at (13.95, -0.35) {
    \textbf{Asymmetric Critic $V_\phi$}\\[0.5pt]
    {\scriptsize Full 157-D state (all conditions)}\\
    {\scriptsize PPO updates $\cdot$ dropped at eval}
  };

  % Vertical observation bus at x = 2.65
  \coordinate (bus_mid)    at (2.65, 1.30);
  \coordinate (bus_top)    at (2.65, 3.15);
  \coordinate (bus_bypass) at (2.65, 2.35);
  \coordinate (bus_ok)     at (2.65, 1.40);
  \coordinate (bus_no)     at (2.65, 0.72);
  \coordinate (bus_bot)    at (2.65, -0.35);

  \draw[line width=0.95pt, draw=navink!80] (env.east) -- (bus_mid)
    node[midway, above, font=\sffamily\scriptsize, text=navink] {$s_t$ (157-D)};
  \draw[line width=0.95pt, draw=navink!80] (bus_bot) -- (bus_top);
  \fill[navink!80] (bus_mid) circle (1.4pt) (bus_bypass) circle (1.4pt)
                   (bus_ok) circle (1.4pt) (bus_no) circle (1.4pt);

  % Top branch: Action-level intent pipeline
  \draw[wire] (bus_top) -- (expert.west)
    node[midway, above, font=\sffamily\scriptsize, text=navink] {157-D};
  \draw[wire, draw=intentblue] (expert.east) -- (pca.west)
    node[midway, above, font=\sffamily\scriptsize, text=intentblue] {20-D $a_t^E$};
  \draw[dashwire] (basis.south) -- (pca.north)
    node[midway, right, font=\sffamily\scriptsize, text=mutedgray] {fixed axes};
  \draw[wire, draw=intentblue] (pca.east) -- (channel.west)
    node[midway, above, font=\sffamily\scriptsize, text=intentblue] {$z_t \in \mathbb{R}^k$};

  % Bypass branch: Object / Goal intents (no expert, no PCA)
  \draw[wire, dashed, draw=mutedgray]
    (bus_bypass) -- (12.85, 2.35) -- ([xshift=-1.1cm]channel.south);
  \node[above, font=\sffamily\scriptsize, text=mutedgray] at (7.75, 2.35)
    {Obj / Goal intents: direct from 157-D $s_t$ (no expert, no PCA)};

  % Channel -> Concat
  \draw[wire, draw=chanorange]
    (channel.south) -- (13.95, 1.93) -- (10.55, 1.93) -- (concat.north);
  \node[above, font=\sffamily\scriptsize, text=chanorange!90!black] at (12.25, 1.93)
    {$\tilde{z}_t$ ($k$-D sole goal cue)};

  % Middle branch: 149-D state -> Concat -> Actor, and 8-D goal stripped
  \draw[wire] (bus_ok) -- (concat.west)
    node[midway, above, font=\sffamily\scriptsize, text=navink]
    {149-D physical state (\texttt{obs[0:64,72:157]}: 60~Hz proprio \& obj)};
  \draw[wire, draw=intentblue] (concat.east) -- (actor.west |- concat.east)
    node[midway, above, font=\sffamily\scriptsize, text=intentblue] {$149+k$};

  \coordinate (strip_end) at (6.7, 0.72);
  \draw[stripwire] (bus_no) -- (strip_end)
    node[midway, below, font=\sffamily\scriptsize, text=stripcrimson] {8-D goal (\texttt{obs[64:72]})};
  \node[circle, draw=stripcrimson, fill=white, line width=0.9pt, inner sep=1pt,
        font=\sffamily\scriptsize\bfseries, text=stripcrimson] (cross) at (strip_end) {$\times$};
  \node[right=2.5pt of cross, font=\sffamily\scriptsize\bfseries, text=stripcrimson]
    {Stripped (goal-hidden)};

  % Bottom branch: Full 157-D -> Critic -> Actor (solid)
  \draw[criticwire] (bus_bot) -- (critic.west)
    node[midway, above, font=\sffamily\scriptsize, text=criticgreen!85!black]
    {Full 157-D state (goal preserved, unmasked)};
  \draw[criticwire] (critic.north) -- (actor.south)
    node[midway, left, font=\sffamily\scriptsize, text=criticgreen!85!black]
    {PPO $\hat{A}_t$};

  % Action feedback loop + Step 5 evaluation note
  \draw[wire, draw=intentblue]
    (actor.east) -- (16.35, 1.25) -- (16.35, -1.12) -- (-0.35, -1.12) -- (env.south);
  \node[fill=white, inner sep=2pt, font=\sffamily\scriptsize, text=intentblue]
    at (8.00, -1.12)
    {20-D Hand action $a_t$ (60~Hz) \quad $\vert$ \quad Step 5 Eval: Actor + Channel only (25 ckpts $\times$ 1,024 eps, no Critic)};

\end{tikzpicture}%
}

  \caption{Part~I student training and channel-sweep pipeline (Steps~1--5). Action intents project the expert's 20-D action onto a fixed PCA basis; the actor concatenates $\tilde{z}_t$ with the 149-D state under a 157-D training critic.}
  \label{fig:protocol}
\end{figure}

At each control step $t$, we compute $z_t \in \mathbb{R}^k$ across three abstraction levels (Table~\ref{tab:intent_taxonomy}):
\begin{enumerate}[leftmargin=1.5em,itemsep=1.5pt,topsep=2pt]
  \item \textbf{Action-level intents:} Using PCA bases fitted on 1,228,800 deterministic mean-action transitions from the frozen reference expert (maximizing retained action variance), \texttt{hand-$k$} ($k \in \{2,4,8,20\}$) projects the expert's online mean action $a_t^E$ onto the top $k$ whole-hand PCs: $z_t = W_{:,1:k}^\top (a_t^E - \mu)$. To test wrist--finger decoupling \citep{liu2016analysis}, \texttt{wrist-2} sends $a^E_{0:2}$ and \texttt{wrist2-finger-$k$} ($k \in \{2,6\}$) appends the top $k$ finger-only PCs. On Task~2 we also test \texttt{joint-$k$}, projecting $a_t^E$ onto the joint-position PCA basis.
  \item \textbf{Object-subgoal intents:} Computed directly from object and goal poses (no reference expert needed): \texttt{obj-1} sends the remaining geodesic angle $\theta_t$; \texttt{obj-4} sends an intermediate waypoint $\Delta\theta = 0.5$~rad ahead along $\mathrm{slerp}(q_{\mathrm{obj}}, q_{\mathrm{goal}}, \min(1, 0.5/\theta_t))$.
  \item \textbf{Goal-level intents:} Also computed directly from the target: \texttt{goal-4} sends canonicalized $q_{\mathrm{goal}}$ ($w \ge 0$); \texttt{goal-yaw} sends $[\sin\psi_{\mathrm{goal}}, \cos\psi_{\mathrm{goal}}]$ (2-D); \texttt{goal-7} sends $[q_{\mathrm{goal}}, \sin\phi_{\mathrm{goal}}, \cos\phi_{\mathrm{goal}}, \mathbb{I}_{\mathrm{turn}}]$ (7-D) for unified Task~3.
\end{enumerate}

\begin{table}[t]
  \centering
  \caption{Summary of evaluated low-dimensional intent interfaces across abstraction levels.}
  \label{tab:intent_taxonomy}
  \small
  \begin{tabular}{llcl}
    \toprule
    Level & Interface Name & Dim $k$ & Definition / Physical Meaning \\
    \midrule
    None & \texttt{blind} & 0 & No intent (goal-hidden lower bound) \\
    \multirow{4}{*}{Action} & \texttt{wrist-2} & 2 & Raw reference-expert wrist commands $a^E_{0:2}$ \\
     & \texttt{hand-$k$} ($k \in \{2,4,8,20\}$) & $k$ & Top-$k$ PCA projection of 20-D expert action $a^E_t$ \\
     & \texttt{wrist2-finger-$k$} ($k \in \{2,6\}$) & $2+k$ & Wrist commands + top-$k$ 18-D finger PCA \\
     & \texttt{joint-$k$} ($k \in \{2,4\}$, Task~2) & $k$ & Expert action projected onto joint-position PCA basis \\
    \multirow{2}{*}{Object Subgoal} & \texttt{obj-1} & 1 & Remaining geodesic rotation angle $\theta_t$ \\
     & \texttt{obj-4} & 4 & Geodesic waypoint $\Delta\theta=0.5$~rad ahead of cube \\
    \multirow{3}{*}{Goal} & \texttt{goal-4} & 4 & Target cube orientation quaternion $q_{\mathrm{goal}}$ \\
     & \texttt{goal-yaw} (Tasks~2--3) & 2 & Target yaw / hinge angle $[\sin\psi, \cos\psi]$ \\
     & \texttt{goal-7} (Task~3 unified) & 7 & $q_{\mathrm{goal}}$ (4) + hinge target $[\sin,\cos]$ (2) + \texttt{is\_turn} (1) \\
    \bottomrule
  \end{tabular}
\end{table}

\subsection{Channel Degradation Models and Waiting-Time Law (Steps 3--4)}
\label{sec:method_channel}
To model a slow, band-limited planner, $z_t$ passes through three controlled degradations:
(i)~\textbf{Zero-order hold refresh:} recomputed every $N \in \{1, 6, 12, 30\}$ steps (60, 10, 5, 2~Hz) and held in between, $\tilde{z}_t = z_{N \lfloor t/N \rfloor}$ (tested both evaluation-only and trained-under-hold), alongside a \textbf{Slow Expert} control whose own 20-D action is held for $N$ steps despite full goal visibility;
(ii)~\textbf{Fixed latency:} $\tilde{z}_t = z_{\max(0, t-d)}$ for $d \in \{6, 30, 60, 120\}$ steps (0.1--2.0~s, evaluation-only and trained-under-delay); and
(iii)~\textbf{Scalar quantization:} uniform $b \in \{8, 4, 3, 2\}$-bit quantization per component on $[-1,1]$ (evaluation-only).

\paragraph{Analytical waiting-time law.}
When a goal-level intent remains valid until goal completion, channel degradation primarily adds dead time $w$ after each completion before the new target arrives ($w=d$ for delay; $w \approx N/2$ for periodic hold).
With zero-wait completion time $T_0$, expected completions per 600-step episode obey:
\begin{equation}
  P(w) \approx \frac{600}{T_0 + w}, \qquad P_{\mathrm{cycle}}(w) \approx \frac{600}{T_0 + L \cdot w},
  \label{eq:waiting}
\end{equation}
where $L$ is the number of sequential phases per cycle ($L=2$ for \texttt{up,turn}; $L=3$ for \texttt{any,up,turn}).

\subsection{Multi-Skill Control and End-to-End VLM System (Steps 6--7)}
\label{sec:method_unified_vlm}
\paragraph{Step~6: Multi-skill 60~Hz controller (\texttt{Unified v4}).}
For Task~3, we compare \textbf{Two-Expert Handoff} (switching between the Stage~1 expert and a face-turn expert fine-tuned on 34,243 Stage~1 \texttt{up} end states) against a single 60~Hz goal-conditioned policy, \textbf{\texttt{Unified v4}} (162-D observation adding hinge error, \texttt{is\_turn}, hold count, and hinge velocity).
Resolving three RL exploits observed in v1--v3 (Appendix~\ref{app:unified_debug}), \texttt{Unified v4} (i)~terminates episodes on 120-step phase timeout with $V(s')=0$, (ii)~initializes 50\% of resets into $\pm 90^\circ$ \texttt{turn} phases from Stage~1 \texttt{up} handoff states, and (iii)~uses monotonic hold-milestone shaping ($+5$/step, cap $+50$) with fall penalty $-15$.

\paragraph{Step~7: Upstream VLM adaptation and online 60~Hz closed loop (Figure~\ref{fig:vlm_arch}).}
Guided by Part~I, we freeze \texttt{Unified v4} as the 60~Hz goal-conditioned execution layer and discretize the hinged cube into a 96-state graph (24 axis-aligned orientations $\times$ 4 detents, diameter $L=4$).
Macro-actions are JSON goals: \texttt{\{"action":"reorient","up":$c_1$,"front":$c_2$\}} or \texttt{\{"action":"turn","delta":$\pm 90$\}} (legal only when \texttt{up=="white"}).
States use a visually observable \texttt{relational} format (listing \texttt{up}, \texttt{front}, and the 1/3 top-strip color aligned above each 2/3 base side), and visual observations tile three fixed cameras (\texttt{c1 + s040 + s220}) into one $1164\times 412$ \texttt{montage} image.
We adapt \texttt{Qwen3-VL-8B-Instruct} via LoRA ($r=16, \alpha=32$, ViT frozen; 0.51\% trainable parameters):
\textbf{\texttt{C1-v1}} (2 epochs) learns structured \emph{Map-then-Reason} outputs (\texttt{view\_obs $\to$ state $\to$ thought $\to$ action});
\textbf{\texttt{C1-v2}} (1 epoch) adds multi-template phrasing, calibrated unanswerable-frame abstention (CBD), 30\% color-permutation augmentation, and explicit cyclic $\mathbb{Z}_4$ turn CoT; and
\textbf{\texttt{C2} GRPO} (\texttt{C2-Terminal} / \texttt{C2-Grounded}) applies multi-turn outcome-reward RL \citep{shao2024deepseekmath} over a 96-state real-rendered image pool before real-time 60~Hz Isaac Lab verification.

\begin{figure}[t]
  \centering
  % Compact TikZ diagram for Figure 2: Part II End-to-End VLM + Frozen Unified v4 System (Steps 6--7)
\definecolor{navink}{HTML}{17213A}
\definecolor{mutedgray}{HTML}{62708D}
\definecolor{cardbg}{HTML}{F6F8FC}
\definecolor{cardborder}{HTML}{CBD5E1}
\definecolor{intentblue}{HTML}{2C67F2}
\definecolor{intentbluebg}{HTML}{EEF4FF}
\definecolor{vlmviolet}{HTML}{6D28D9}
\definecolor{vlmvioletbg}{HTML}{F5F3FF}
\definecolor{ctrlgreen}{HTML}{1E9C77}
\definecolor{ctrlgreenbg}{HTML}{EDF9F5}

\resizebox{\linewidth}{!}{%
\begin{tikzpicture}[
    font=\sffamily\footnotesize,
    >=Stealth,
    box/.style={
      rectangle, rounded corners=3.5pt, draw=cardborder, fill=cardbg,
      line width=0.8pt, align=center, inner sep=4.5pt, text=navink
    },
    envbox/.style={
      box, draw=navink!65, fill=cardbg, line width=1.0pt,
      minimum width=3.3cm, minimum height=1.45cm
    },
    vlmbox/.style={
      box, draw=vlmviolet, fill=vlmvioletbg, line width=1.0pt,
      minimum width=5.2cm, minimum height=1.45cm
    },
    ctrlbox/.style={
      box, draw=ctrlgreen, fill=ctrlgreenbg, line width=1.0pt,
      minimum width=4.1cm, minimum height=1.45cm
    },
    wire/.style={->, line width=0.95pt, draw=navink!80}
  ]

  % Three main blocks with generous arrow gaps so labels never overlap boxes
  \node[envbox] (env) at (0, 0) {
    \textbf{Isaac Lab 60~Hz Env}\\[1.5pt]
    {\scriptsize Hinged Cube (96 states)}\\
    {\scriptsize 3 fixed cams: \texttt{c1, s040, s220}}
  };

  \node[vlmbox] (vlm) at (7.45, 0) {
    \textbf{Slow VLM Planner (Step 7)}\\[1.5pt]
    {\scriptsize \texttt{Qwen3-VL-8B} (Frozen ViT + 0.51\% LoRA)}\\
    {\scriptsize Map-then-Reason: \texttt{view\_obs} $\to$ \texttt{state} $\to$ CoT $\to$ JSON}
  };

  \node[ctrlbox] (ctrl) at (16.15, 0) {
    \textbf{Frozen 60~Hz Controller}\\[1.5pt]
    {\scriptsize \texttt{Unified v4} (Step 6, goal-conditioned)}\\
    {\scriptsize $\sim$16 steps/macro-action $\cdot$ 0\% drop (4~s hold)}
  };

  % Top band: Slow Macro-Loop
  \draw[wire, draw=vlmviolet] (env.east) -- (vlm.west)
    node[midway, above, font=\sffamily\scriptsize, text=vlmviolet, align=center] {3-cam \texttt{montage}\\{\tiny + target text}};

  \draw[wire, draw=intentblue] (vlm.east) -- (ctrl.west)
    node[midway, above, font=\sffamily\scriptsize, text=intentblue, align=center] {Goal macro-action\\{\tiny \texttt{reorient} / \texttt{turn}($\pm 90^\circ$)}};

  % Bottom band: Fast 60 Hz Motor Closed Loop (Ctrl <-> Env)
  \draw[wire, draw=ctrlgreen]
    (ctrl.south) -- (16.15, -1.25) -- (0, -1.25) -- (env.south);
  \node[fill=white, inner sep=2pt, font=\sffamily\scriptsize, text=ctrlgreen!90!black]
    at (8.05, -1.25) {\textbf{Fast 60~Hz Motor Closed Loop:} 20-D Shadow Hand actions $a_t$ from proprio + cube state ($s_t$) until goal reached};

  % Subtle loop-rate annotations above
  \node[font=\sffamily\scriptsize\bfseries, text=vlmviolet] at (7.45, 1.05)
    {Slow Deliberative Loop ($\sim$0.2--2~Hz, goal-level commands only)};

\end{tikzpicture}%
}

  \caption{Part~II end-to-end VLM manipulation architecture (Steps~6--7). Frozen \texttt{Unified v4} acts as the 60~Hz goal-conditioned skill controller executing discrete macro-actions emitted by the slow Map-then-Reason VLM from live 3-camera \texttt{montage} images.}
  \label{fig:vlm_arch}
\end{figure}
